% !TEX root = ../main.tex%

\chapter{The Transfinite Phoa Principle and the universal \texorpdfstring{$L$}{L}-(co)algebra}
\label{chapter-omega}

This chapter gives a concrete construction of the initial $L$-algebra and the final $L$-coalgebra, both of which will play a pivotal role in investigating the general recursion in SDT. The construction is based on the relation between the lifting functor and the simplices revealed in Chapter \ref{chapter-interval}. Furthermore, this construction will lead to a generalisation of the Phoa principle for the initial $L$-algebra, which enables us to compare the topology of the transfinite counterpart of the simplex and the spine.

\section{The Initial \texorpdfstring{$L$}{L}-algebra and the Final \texorpdfstring{$L$}{L}-coalgebra}
We introduce the concept of the initial $L$-algebra and the final $L$-coalgebra in the language of category theory.
\begin{definition}
    Let $\alpha:LA\to A$ and $\beta:LB\to B$ be two $L$-algebras. A function $f:A\to B$ is a \textbf{homomorphism between the $L$-algebras} $\alpha$ and $\beta$ iff the following diagram commutes:
    % https://q.uiver.app/#q=WzAsNCxbMCwwLCJMQSJdLFswLDEsIkEiXSxbMSwwLCJMQiJdLFsxLDEsIkIiXSxbMSwzLCJmIiwyXSxbMCwyLCJMZiJdLFswLDEsIlxcYWxwaGEiLDJdLFsyLDMsIlxcYmV0YSJdXQ==
    \[\begin{tikzcd}
        LA & LB \\
        A & B
        \arrow["Lf", from=1-1, to=1-2]
        \arrow["\alpha"', from=1-1, to=2-1]
        \arrow["\beta", from=1-2, to=2-2]
        \arrow["f"', from=2-1, to=2-2]
    \end{tikzcd}\]
    The \textbf{category of $L$-algebras} is a category where objects are $L$-algebras and the morphisms are homomorphisms between the $L$-algebras. Dually, we also have the \textbf{category of $L$-coalgebras}.
\end{definition}
\begin{definition}
    The \textbf{initial $L$-algebra} is the initial object of the category of $L$-algebras, denoted $\phi:L\w\to\w$. The \textbf{final $L$-coalgebra} is the final object of the category of $L$-coalgebras, denoted $\psi:\wb\to L\wb$.
\end{definition}
% \begin{theorem}
%     The initial $L$-algebra and the final $L$-coalgebra exist.
% \end{theorem}

One may attempt to construct the initial $L$-algebra and the final $L$-coalgebra using the Ad\'amek's fixed point theorem \cite{adamek_free_1974}, which asserts that the initial $L$-algebra is the colimit of the following diagram:
% https://q.uiver.app/#q=WzAsNSxbMCwwLCJcXGJvdCJdLFsxLDAsIkxcXGJvdCJdLFsyLDAsIkxeMlxcYm90Il0sWzMsMCwiTF4zXFxib3QiXSxbNCwwLCJcXGNkb3RzIl0sWzAsMSwiISJdLFsxLDIsIkwhIl0sWzIsMywiTF4yISJdLFszLDRdXQ==
\[\begin{tikzcd}
    \bot & {L\bot} & {L^2\bot} & {L^3\bot} & \cdots
    \arrow["{!}", from=1-1, to=1-2]
    \arrow["{L!}", from=1-2, to=1-3]
    \arrow["{L^2!}", from=1-3, to=1-4]
    \arrow[from=1-4, to=1-5]
\end{tikzcd}\]
Unfortunately, this construction only works if $L$ preserves the $\omega$-chains, or equivalently, the initial $L$-algebra is inductive, which was once believed to be true \cite{hyland_first_1991}, but then disproved with counterexamples in the effective topos \cite{van_oosten_axioms_2000}. To proceed with our investigation, we treat this as an axiom.
\begin{axiom}
    \label{ax-inductive}
    The initial $L$-algebra $\omega$ exists and satisfies $\omega=\varinjlim_{n\in\N}L^n\bot$.
\end{axiom}

It is difficult to illustrate this colimit directly. However, Theorem \ref{thm-l-delta} gives rise to an isomorphic chain built from simplices:
% https://q.uiver.app/#q=WzAsMTAsWzAsMCwiXFxib3QiXSxbMSwwLCJMXFxib3QiXSxbMiwwLCJMXjJcXGJvdCJdLFszLDAsIkxeM1xcYm90Il0sWzQsMCwiXFxjZG90cyJdLFswLDEsIlxcRGVsdGFfey0xfSJdLFsxLDEsIlxcRGVsdGFfMCJdLFsyLDEsIlxcRGVsdGFfMSJdLFszLDEsIlxcRGVsdGFfMiJdLFs0LDEsIlxcY2RvdHMiXSxbMCwxLCIhIl0sWzEsMiwiTCEiXSxbMiwzLCJMXjIhIl0sWzMsNF0sWzUsNiwic18wIl0sWzYsNywic18xIl0sWzcsOCwic18yIl0sWzgsOV0sWzAsNSwiXFxjb25nIl0sWzEsNiwiXFxjb25nIl0sWzIsNywiXFxjb25nIl0sWzMsOCwiXFxjb25nIl1d
\[\begin{tikzcd}
    \bot & {L\bot} & {L^2\bot} & {L^3\bot} & \cdots \\
    {\Delta^{-1}} & {\Delta^0} & {\Delta^1} & {\Delta^2} & \cdots
    \arrow["{!}", from=1-1, to=1-2]
    \arrow["\cong", from=1-1, to=2-1]
    \arrow["{L!}", from=1-2, to=1-3]
    \arrow["\cong", from=1-2, to=2-2]
    \arrow["{L^2!}", from=1-3, to=1-4]
    \arrow["\cong", from=1-3, to=2-3]
    \arrow[from=1-4, to=1-5]
    \arrow["\cong", from=1-4, to=2-4]
    \arrow["{s_0}", from=2-1, to=2-2]
    \arrow["{s_1}", from=2-2, to=2-3]
    \arrow["{s_2}", from=2-3, to=2-4]
    \arrow[from=2-4, to=2-5]
\end{tikzcd}\]
where $!=\rec_\bot$. We only need to find the exact form of $s_n:\Delta^{n-1}\to\Delta^n$. Recall that $L\bot=\set{(0,\rec_\bot)}$ and $Lf=\lambda(i,x).(i,f\circ x)$. Thus,
$$L!(0,\rec_\bot)=(0,\rec_\bot\circ\rec_\bot)=(0,(0,\rec_\bot))$$
which corresponds to $0\in\I=\Delta^1$. This means that $s_1(*)=0$ and more generally
$$s_n(i_1,\dots,i_{n-1})=(i_1,\dots,i_{n-1},0)$$
Therefore,
$$\Delta^\omega:=\varinjlim_{n\in\N}s_n=\prn{\sum_{n\in\N}\Delta^n}/(x\sim s_n(x))$$
We then further investigate what the elements of $\Delta^\omega$ exactly are. From the definition, $\Delta^\omega$ should contain all descending sequences of $\I$ of any finite length, up to injections by $s_n$. Let
$$\Delta^\infty:=\set{i\in\I^\N:\forall n\in\N.i(n)\sqsupseteq i(n+1)}$$
We can see that any $\Delta^n$ embeds into $\Delta^\infty$ via
$$\iota_n(i_1,\dots,i_n):=(i_1,\dots,i_n,0,\dots)$$
and that embedding commutes to $s_n$. Hence, $\Delta^\omega$ is exactly the subset of $\Delta^\infty$ containing the images of the embeddings. That is,
$$\Delta^\omega=\set{i\in\Delta^\infty:\exists n\in\N.i(n)=0}$$

We now construct the final $L$-coalgebra, which can be done by calculating the limit of the diagram
% https://q.uiver.app/#q=WzAsMTAsWzAsMCwiXFx0b3AiXSxbMSwwLCJMXFx0b3AiXSxbMiwwLCJMXjJcXHRvcCJdLFszLDAsIkxeM1xcdG9wIl0sWzQsMCwiXFxjZG90cyJdLFswLDEsIlxcRGVsdGFfMCJdLFsxLDEsIlxcRGVsdGFfMSJdLFsyLDEsIlxcRGVsdGFfMiJdLFszLDEsIlxcRGVsdGFfMyJdLFs0LDEsIlxcY2RvdHMiXSxbMiwxLCJMKiIsMl0sWzMsMiwiTF4yKiIsMl0sWzQsM10sWzYsNSwiZF8wIiwyXSxbNyw2LCJkXzEiLDJdLFs4LDcsImRfMiIsMl0sWzksOF0sWzAsNSwiXFxjb25nIl0sWzEsNiwiXFxjb25nIl0sWzIsNywiXFxjb25nIl0sWzMsOCwiXFxjb25nIl0sWzEsMCwiKiIsMl1d
\[\begin{tikzcd}
    \top & {L\top} & {L^2\top} & {L^3\top} & \cdots \\
    {\Delta^0} & {\Delta^1} & {\Delta^2} & {\Delta^3} & \cdots
    \arrow["\cong", from=1-1, to=2-1]
    \arrow["{*}"', from=1-2, to=1-1]
    \arrow["\cong", from=1-2, to=2-2]
    \arrow["{L*}"', from=1-3, to=1-2]
    \arrow["\cong", from=1-3, to=2-3]
    \arrow["{L^2*}"', from=1-4, to=1-3]
    \arrow["\cong", from=1-4, to=2-4]
    \arrow[from=1-5, to=1-4]
    \arrow["{d_0}"', from=2-2, to=2-1]
    \arrow["{d_1}"', from=2-3, to=2-2]
    \arrow["{d_2}"', from=2-4, to=2-3]
    \arrow[from=2-5, to=2-4]
\end{tikzcd}\]
Similarly to the previous observations, it can be seen that $d_1(i,j)=(i)$, and more generally,
$$d_n(i_1,\dots,i_n)=(i_1,\dots,i_{n-1})$$
One can immediately realise that
$$\varprojlim_{n\in\N}d_n=\Delta^\infty$$
\begin{definition}
    $\Delta^\omega$ is the $\omega$-simplex and $\Delta^\infty$ is the $\infty$-simplex.
\end{definition}
\begin{theorem}
    $\Delta^\omega\cong\w$ carries the initial $L$-algebra
    $$\phi(i,j)=(i,\sigma_\omega(i,j))$$
    where
    $$\sigma_\omega(i,j:\sem{i}\to\I^\N):=\lambda n.\sigma(i,j(-)(n))$$

    $\Delta^\infty\cong\wb$ carries the final $L$-coalgebra
    $$\psi(i)=(i(0),\lambda\_.\lambda n.i(n+1))$$
\end{theorem}

Similarly to $\Delta_n$, we also have
$$\Delta_\infty:=\set{i\in\I^\N:\forall n\in\N.i(n)\sqsubseteq i(n+1)}$$
and
$$\Delta_\omega=\set{i\in\Delta_\infty:\exists n\in\N.i(n)=1}$$
Though these objects do not occur as an initial(final) $L$-coalgebra, we will find their connection with $\Delta^\omega$ and $\Delta^\infty$ via a transfinite variant of the Phoa principle.

\section{The Transfinite Phoa Principle}
Recall that the higher Phoa principle samples a function $\Delta^n\to\I$ on the $n+1$ vertices. We now repeat this sampling process on $\Delta^\omega$. It is not difficult to see that the vertices mapped from $\Delta^n$ to $\Delta^\omega$ via $\iota_n$ are
$$v_k:=(1^k\overline{0})=(\overbrace{1,\dots,1}^\text{$k$ copies},0,\dots)$$
which defines an embedding $\N\hookrightarrow\Delta^\omega$.
\begin{lemma}
    For any function $f:\Delta^\omega\to\I$, the evaluation
    $$j:=\lambda n.f(v_n)=f(1^n\overline{0})$$
    gives an increasing sequence, and hence $j\in\Delta_\infty$.
\end{lemma}
\begin{proof}
    For any $n\ge1$, $v_{n-1},v_n\in\Delta^\omega$ can be seen as embedded from the corresponding vertices in the $n$-simplex. By precomposing the embedding $\Delta^n\hookrightarrow\Delta^\omega$ with $f$, we get its restriction on $\Delta^n$, which is monotonic by Lemma \ref{lemma-mono-n}. Therefore, $f(v_{n-1})\sqsubseteq f(n)$.
\end{proof}

This observation is enough for us to conjecture that $\I^{\Delta^\omega}\cong\Delta_\infty$, which we call the transfinite Phoa principle. Since both $\Delta^\omega$ and $\Delta_\infty$ are constructed by (co)limit, we can prove the transfinite Phoa principle from the higher version.
\begin{lemma}
    \label{lemma-chain}
    The higher Phoa principle gives an isomorphism between chains
    % https://q.uiver.app/#q=WzAsMTAsWzAsMCwiXFxtYXRoYmJ7SX1ee1xcRGVsdGFfey0xfX0iXSxbMSwwLCJcXG1hdGhiYntJfV57XFxEZWx0YV8wfSJdLFsyLDAsIlxcbWF0aGJie0l9XntcXERlbHRhXzF9Il0sWzMsMCwiXFxtYXRoYmJ7SX1ee1xcRGVsdGFfMn0iXSxbNCwwLCJcXGNkb3RzIl0sWzAsMSwiXFxEZWx0YV4wIl0sWzEsMSwiXFxEZWx0YV4xIl0sWzIsMSwiXFxEZWx0YV4yIl0sWzMsMSwiXFxEZWx0YV4zIl0sWzQsMSwiXFxjZG90cyJdLFsyLDEsIlxcbWF0aGJie0l9XntzXzF9IiwyXSxbMywyLCJcXG1hdGhiYntJfV57c18yfSIsMl0sWzQsM10sWzYsNSwiZF4wIiwyXSxbNyw2LCJkXjEiLDJdLFs4LDcsImReMiIsMl0sWzksOF0sWzAsNSwiXFxjb25nIl0sWzEsNiwiXFxjb25nIl0sWzIsNywiXFxjb25nIl0sWzMsOCwiXFxjb25nIl0sWzEsMCwiXFxtYXRoYmJ7SX1ee3NfMH0iLDJdXQ==
    \[\begin{tikzcd}
        {\mathbb{I}^{\Delta^{-1}}} & {\mathbb{I}^{\Delta^0}} & {\mathbb{I}^{\Delta^1}} & {\mathbb{I}^{\Delta^2}} & \cdots \\
        {\Delta_0} & {\Delta_1} & {\Delta_2} & {\Delta_3} & \cdots
        \arrow["\cong", from=1-1, to=2-1]
        \arrow["{\mathbb{I}^{s_0}}"', from=1-2, to=1-1]
        \arrow["\cong", from=1-2, to=2-2]
        \arrow["{\mathbb{I}^{s_1}}"', from=1-3, to=1-2]
        \arrow["\cong", from=1-3, to=2-3]
        \arrow["{\mathbb{I}^{s_2}}"', from=1-4, to=1-3]
        \arrow["\cong", from=1-4, to=2-4]
        \arrow[from=1-5, to=1-4]
        \arrow["{d^0}"', from=2-2, to=2-1]
        \arrow["{d^1}"', from=2-3, to=2-2]
        \arrow["{d^2}"', from=2-4, to=2-3]
        \arrow[from=2-5, to=2-4]
    \end{tikzcd}\]
    where
    $$d^n:\Delta_{n+1}\to\Delta_n:=\lambda(i_1,\dots,i_{n+1}).(i_1,\dots,i_n)$$
\end{lemma}
\begin{proof}
    Let $f:\Delta^n\to\I$ and $j=(j_0,\dots,j_n):=(f(v_0),\dots,f(v_n))\in\Delta_{n+1}$. For any $(i_1,\dots,i_{n-1})\in\Delta^{n-1}$, we have
    $$\begin{aligned}
        &\I^{s_n}(f)(i_1,\dots,i_{n-1}) \\
        =&(f\circ s_n)(i_1,\dots,i_{n-1}) \\
        =&f(s_n(i_1,\dots,i_{n-1})) \\
        =&f(i_1,\dots,i_{n-1},0) \\
    \end{aligned}$$
    and hence $\I^{s_n}(f)(v_k)=f(v_k)$ for $v_0,\dots,v_{n-1}\in\Delta^{n-1}$. Thus,
    $$(\I^{s_n}(f)(v_0),\dots,\I^{s_n}(f)(v_{n-1}))=(j_0,\dots,j_{n-1})=d^n(j)$$
    Therefore, the diagram between the chains commutes, and the higher Phoa principle gives an isomorphism between them.
\end{proof}
\begin{theorem}[Transfinite Phoa principle, \texttt{Omega.Phoaω}]
    \label{thm-phoa-omega}
    The evaluation on vertices $v:\N\hookrightarrow\Delta^\omega$ gives an isomorphism
    $$\I^{\Delta^\omega}\cong\Delta_\infty$$
\end{theorem}
\begin{proof}
    From Lemma \ref{lemma-chain} and the property of the internal-hom, we have
    $$\I^{\Delta^\omega}=[\I,\Delta^\omega]=[\I,\varinjlim_{n\in\N}s_n]=\varprojlim_{n\in\N}[\I,s_n]=\varprojlim_{n\in\N}\I^{s_n}\cong\varprojlim_{n\in\N}d^n=\Delta_\infty$$
\end{proof}
One can realise that this proof by (co)limits only works for $\I^{\Delta^n}\cong\Delta_{n+1}$, and fails to proceed with $\I^{\Delta^n}\cong\Delta^{n+1}$ since the diagram no longer commutes if we replace $d^n$ with $d_n$. This makes a stronger argument for why we should treat $\Delta^n$ and $\Delta_n$ as dual but different objects, even though we obviously have $\Delta^n\cong\Delta_n$ for the finite case.

\section{The \texorpdfstring{$\omega$}{omega}-spine}
Similarly to the simplices, there also exists a transfinite counterpart for the spines, denoted $\Lambda_\omega$. Recall that a $n$-spine $\Lambda_n$ contains the minimal information to describe the order on $[n]=\set{0,\dots,n}$. A transfinite version of the spine should likewise grasp the information that generates the order on $\N$, which we can see is the set of relations $\set{n\le n+1:n\in\N}$. By extending the diagram used to define $\Lambda_n$, we can define $\Lambda_\omega$ as the limit of the following diagram:
% https://q.uiver.app/#q=WzAsNyxbMCwwLCIwIl0sWzIsMCwiMSJdLFs0LDAsIjIiXSxbMSwxLCIwXFxsZTEiXSxbMywxLCIxXFxsZTIiXSxbNSwxLCJcXGNkb3RzIl0sWzUsMCwiXFxjZG90cyJdLFswLDMsIiIsMCx7InN0eWxlIjp7InRhaWwiOnsibmFtZSI6Imhvb2siLCJzaWRlIjoidG9wIn19fV0sWzEsMywiIiwyLHsic3R5bGUiOnsidGFpbCI6eyJuYW1lIjoiaG9vayIsInNpZGUiOiJ0b3AifX19XSxbMSw0LCIiLDAseyJzdHlsZSI6eyJ0YWlsIjp7Im5hbWUiOiJob29rIiwic2lkZSI6InRvcCJ9fX1dLFsyLDQsIiIsMix7InN0eWxlIjp7InRhaWwiOnsibmFtZSI6Imhvb2siLCJzaWRlIjoidG9wIn19fV0sWzIsNSwiIiwwLHsic3R5bGUiOnsidGFpbCI6eyJuYW1lIjoiaG9vayIsInNpZGUiOiJ0b3AifX19XV0=
\[\begin{tikzcd}
    0 && 1 && 2 & \cdots \\
    & {0\le1} && {1\le2} && \cdots
    \arrow[hook, from=1-1, to=2-2]
    \arrow[hook, from=1-3, to=2-2]
    \arrow[hook, from=1-3, to=2-4]
    \arrow[hook, from=1-5, to=2-4]
    \arrow[hook, from=1-5, to=2-6]
\end{tikzcd}\]
Though this diagram is infinite, we can find an equivalent finite limit by indexing the objects in the diagram using $\N$.
\begin{definition}
    The \textbf{$\omega$-spine} $\Lambda_\omega$ is the limit of the following diagram:
    % https://q.uiver.app/#q=WzAsMixbMCwwLCJcXE4iXSxbMCwxLCJcXE5cXHRpbWVzXFxJIl0sWzAsMSwicF8wIiwyLHsiY3VydmUiOjF9XSxbMCwxLCJwXzEiLDAseyJjdXJ2ZSI6LTF9XV0=
    \[\begin{tikzcd}
        \N \\
        {\N\times\I}
        \arrow["{p_0}"', curve={height=6pt}, from=1-1, to=2-1]
        \arrow["{p_1}", curve={height=-6pt}, from=1-1, to=2-1]
    \end{tikzcd}\]
    where $p_0(n):=(n,0)$ and $p_1(n):=(n+1,1)$. Equivalently, $\Lambda_\omega$ is the coequalizer of $p_0$ and $p_1$.
\end{definition}
The $\omega$-spine described as an HIT in directed homotopy is
\begin{code}
data Λω : Type where
    step : ℕ → Λω
    n≤n+1 : ∀ n → step n ⊑ step (suc n)
\end{code}
which can be translated into an HIT in Cubical Agda as
\begin{code}
data Λω : Type where
    step : ℕ → Λω
    n≤n+1 : ℕ → S → Λω
    n≤n+1-s0 : ∀ n → n≤n+1 n s0 ≡ step n
    n≤n+1-s1 : ∀ n → n≤n+1 n s1 ≡ step (suc n)
\end{code}
We can then get the following results:
\begin{theorem}[\texttt{Omega.PhoaΛω}]
    \label{thm-spine-omega}
    The evaluation on vertices $v:\N\hookrightarrow\Lambda_\omega$ gives an isomorphism
    $$\I^{\Lambda_\omega}\cong\Delta_\infty$$
\end{theorem}
\begin{theorem}
    \label{thm-sobrio-omega}
    $\Lambda_\omega$ and $\Delta^\omega$ are sobriomorphic via the isomorphism given by the transfinite Phoa principle.
\end{theorem}
